% GENERATED FILE. Edit canonical lessons, then run npm run generate:books.
% canonical-source-hash: fnv1a64:ccdddaeab622e275
% canonical-lessons: GU-C258-kriket, GU-C258-dado, GU-C258-bet, GU-C258-maarvu, GU-C258-uchalvu

\chapter{Bat and ball}
\label{ch:gu-bat-and-ball}
\hlchaptermodality{258}

\begin{chapteropening}
\textbf{By the end of this chapter:} \emph{I can say who hit the ball, how it bounced and what we played.}

Complete the last lesson of chapter 258: I can say who hit the ball, how it bounced and what we played.
\end{chapteropening}

\section[\texorpdfstring{\gu{ક્રિકેટ} (krikeṭ)}{ક્રિકેટ (krikeṭ)}]{\gu{ક્રિકેટ} (krikeṭ) — cricket}

\label{lesson:GU-C258-kriket}



\begin{quote}
Before the new one: say the Gujarati for to take part, then the Gujarati for to join.
\end{quote}

\subsection*{You'll want to know: \gu{ક્રિકેટ}}
\textbf{\gu{ક્રિકેટ}} — \emph{krikeṭ} — \textquotedblleft{}cricket\textquotedblright{}.

\begin{quote}
\textbf{\gu{અમે} \gu{રવિવારે} \gu{ક્રિકેટ} \gu{રમ્યા}.} — \emph{ame ravivāre krikeṭ ramyā} — \textquotedblleft{}We played cricket on Sunday.\textquotedblright{}
\end{quote}

\begin{cousinweb}
\textbf{\gu{ક્રિકેટ}} (\emph{krikeṭ}) is a loan from English \emph{cricket}, the game.
\end{cousinweb}

\begin{grammarlens}[title={What you've built}]
One more word for everyday life at A2.
\end{grammarlens}

\subsection*{Your turn}
\begin{itemize}
\raggedright
  \item \emph{Say it:} \emph{krikeṭ}
  \item \emph{Say it:} \emph{krikeṭ}, once more
  \item \emph{Say it:} \emph{ame ravivāre krikeṭ ramyā} — We played cricket on Sunday.
  \item \emph{Recall:} say the Gujarati for to take part, then the Gujarati for to join, then say \emph{krikeṭ} again
\end{itemize}

\begin{tcolorbox}[breakable,colback=teal!4,colframe=teal!35!black,title={Before you move on}]
What does \gu{ક્રિકેટ} mean? (\textquotedblleft{}Cricket\textquotedblright{}.) Say it once more.
\end{tcolorbox}
\section[\texorpdfstring{\gu{દડો} (daḍo)}{દડો (daḍo)}]{\gu{દડો} (daḍo) — ball}

\label{lesson:GU-C258-dado}



\begin{quote}
Before the new one: say the Gujarati for to join, then the Gujarati for cricket.
\end{quote}

\subsection*{You'll want to know: \gu{દડો}}
\textbf{\gu{દડો}} — \emph{daḍo} — \textquotedblleft{}ball\textquotedblright{}.

\begin{quote}
\textbf{\gu{છોકરાએ} \gu{દડો} \gu{ફેંક્યો}.} — \emph{chokrāe daḍo phẽkyo} — \textquotedblleft{}The boy threw the ball.\textquotedblright{}
\end{quote}

\begin{cousinweb}
\textbf{\gu{દડો}} (\emph{daḍo}) is inherited. EtymDB traces it to Old Gujarati \emph{daḍau}, \textquotedblleft{}ball\textquotedblright{}, and CDIAL lists Old Gujarati \emph{daïḍaü} and Gujarati \emph{deṛo}, \textquotedblleft{}leather ball\textquotedblright{}, under Sanskrit \emph{dṛ́ti}, \textquotedblleft{}leather bag\textquotedblright{} (CDIAL 6511).
\end{cousinweb}

\begin{grammarlens}[title={What you've built}]
One more word for everyday life at A2.
\end{grammarlens}

\subsection*{Your turn}
\begin{itemize}
\raggedright
  \item \emph{Say it:} \emph{daḍo}
  \item \emph{Say it:} \emph{daḍo}, once more
  \item \emph{Say it:} \emph{chokrāe daḍo phẽkyo} — The boy threw the ball.
  \item \emph{Recall:} say the Gujarati for to join, then the Gujarati for cricket, then say \emph{daḍo} again
\end{itemize}

\begin{tcolorbox}[breakable,colback=teal!4,colframe=teal!35!black,title={Before you move on}]
What does \gu{દડો} mean? (\textquotedblleft{}Ball\textquotedblright{}.) Say it once more.
\end{tcolorbox}
\section[\texorpdfstring{\gu{બેટ} (beṭ)}{બેટ (beṭ)}]{\gu{બેટ} (beṭ) — bat}

\label{lesson:GU-C258-bet}



\begin{quote}
Before the new one: say the Gujarati for cricket, then the Gujarati for ball.
\end{quote}

\subsection*{You'll want to know: \gu{બેટ}}
\textbf{\gu{બેટ}} — \emph{beṭ} — \textquotedblleft{}bat\textquotedblright{}.

\begin{quote}
\textbf{\gu{તેણે} \gu{નવું} \gu{બેટ} \gu{ખરીદ્યું}.} — \emph{teṇe navũ beṭ kharīdyũ} — \textquotedblleft{}He bought a new bat.\textquotedblright{}
\end{quote}

\begin{cousinweb}
\textbf{\gu{બેટ}} (\emph{beṭ}) is a loan from English \emph{bat}, the cricket bat.
\end{cousinweb}

\begin{grammarlens}[title={What you've built}]
One more word for everyday life at A2.
\end{grammarlens}

\subsection*{Your turn}
\begin{itemize}
\raggedright
  \item \emph{Say it:} \emph{beṭ}
  \item \emph{Say it:} \emph{beṭ}, once more
  \item \emph{Say it:} \emph{teṇe navũ beṭ kharīdyũ} — He bought a new bat.
  \item \emph{Recall:} say the Gujarati for cricket, then the Gujarati for ball, then say \emph{beṭ} again
\end{itemize}

\begin{tcolorbox}[breakable,colback=teal!4,colframe=teal!35!black,title={Before you move on}]
What does \gu{બેટ} mean? (\textquotedblleft{}Bat\textquotedblright{}.) Say it once more.
\end{tcolorbox}
\section[\texorpdfstring{\gu{મારવું} (mārvũ)}{મારવું (mārvũ)}]{\gu{મારવું} (mārvũ) — to hit, to strike}

\label{lesson:GU-C258-maarvu}



\begin{quote}
Before the new one: say the Gujarati for ball, then the Gujarati for bat.
\end{quote}

\subsection*{You'll want to know: \gu{મારવું}}
\textbf{\gu{મારવું}} — \emph{mārvũ} — \textquotedblleft{}to hit, to strike\textquotedblright{}.

\begin{quote}
\textbf{\gu{તેણે} \gu{દડો} \gu{દૂર} \gu{માર્યો}.} — \emph{teṇe daḍo dūr māryo} — \textquotedblleft{}He hit the ball far.\textquotedblright{}
\end{quote}

\begin{cousinweb}
\textbf{\gu{મારવું}} (\emph{mārvũ}) is a tadbhava. It comes from Sanskrit \emph{māráyati}, \textquotedblleft{}kills\textquotedblright{}; CDIAL glosses the Gujarati verb \textquotedblleft{}to strike, kill\textquotedblright{} (CDIAL 10066).
\end{cousinweb}

\begin{grammarlens}[title={What you've built}]
One more word for everyday life at A2.
\end{grammarlens}

\subsection*{Your turn}
\begin{itemize}
\raggedright
  \item \emph{Say it:} \emph{mārvũ}
  \item \emph{Say it:} \emph{mārvũ}, once more
  \item \emph{Say it:} \emph{teṇe daḍo dūr māryo} — He hit the ball far.
  \item \emph{Recall:} say the Gujarati for ball, then the Gujarati for bat, then say \emph{mārvũ} again
\end{itemize}

\begin{tcolorbox}[breakable,colback=teal!4,colframe=teal!35!black,title={Before you move on}]
What does \gu{મારવું} mean? (\textquotedblleft{}To hit, to strike\textquotedblright{}.) Say it once more.
\end{tcolorbox}
\section[\texorpdfstring{\gu{ઊછળવું} (ūchaḷvũ)}{ઊછળવું (ūchaḷvũ)}]{\gu{ઊછળવું} (ūchaḷvũ) — to bounce, to leap up}

\label{lesson:GU-C258-uchalvu}



\begin{quote}
Before the new one: say the Gujarati for bat, then the Gujarati for to hit.
\end{quote}

\subsection*{You'll want to know: \gu{ઊછળવું}}
\textbf{\gu{ઊછળવું}} — \emph{ūchaḷvũ} — \textquotedblleft{}to bounce, to leap up\textquotedblright{}.

\begin{quote}
\textbf{\gu{દડો} \gu{દીવાલથી} \gu{ઊછળ્યો}.} — \emph{daḍo dīvālthī ūchaḷyo} — \textquotedblleft{}The ball bounced off the wall.\textquotedblright{}
\end{quote}

\begin{cousinweb}
\textbf{\gu{ઊછળવું}} (\emph{ūchaḷvũ}) is a tadbhava. It comes from \emph{*ut-śalati}, \textquotedblleft{}springs up\textquotedblright{}, a reconstructed form beside Sanskrit \emph{ucchalati}; CDIAL writes the Gujarati verb \emph{uchaḷvũ}, \textquotedblleft{}to leap\textquotedblright{} (CDIAL 1843).
\end{cousinweb}

\begin{grammarlens}[title={What you've built}]
One more word for everyday life at A2.
\end{grammarlens}

\subsection*{Your turn}
\begin{itemize}
\raggedright
  \item \emph{Say it:} \emph{ūchaḷvũ}
  \item \emph{Say it:} \emph{ūchaḷvũ}, once more
  \item \emph{Say it:} \emph{daḍo dīvālthī ūchaḷyo} — The ball bounced off the wall.
  \item \emph{Recall:} say the Gujarati for bat, then the Gujarati for to hit, then say \emph{ūchaḷvũ} again
\end{itemize}

\begin{tcolorbox}[breakable,colback=teal!4,colframe=teal!35!black,title={Before you move on}]
What does \gu{ઊછળવું} mean? (\textquotedblleft{}To bounce, to leap up\textquotedblright{}.) Say it once more.
\end{tcolorbox}
